\documentclass[10pt,a4paper]{amsart}
\usepackage[margin=19mm]{geometry}
\usepackage{amsmath,amssymb,mathtools}
\usepackage[scr=rsfs,cal=euler]{mathalfa}
\usepackage{xfrac}
\usepackage[unicode,hidelinks,bookmarks=false]{hyperref}
\newcommand{\ie}{i.e.\ }
\newcommand{\cf}{cf.\ }
\newcommand{\resp}{respectively\ }
\newcommand{\Subsection}{\subsection}
\providecommand{\nobreakdash}{\nobreak\hbox{-}}
\setlength{\emergencystretch}{2em}
\multlinegap0pt
\usepackage{amssymb}
\usepackage{mathtools}
\usepackage{xcolor}
\usepackage{tikz}
\usepackage{float}
\newtheorem{lemma}{Lemma}
\newtheorem{theorem}{Theorem}
\DeclareMathOperator{\supp}{supp}
\DeclareMathOperator{\topo}{top}
\title{Invariant measures with full support and approximation by zero-entropy systems in the $C^0$-Gromov--Hausdorff topology}
\author{Jorge Crisóstomo, Richard Cubas}
\date{Source: arXiv:2604.02810v1 (CC BY 4.0)}
\begin{document}
\maketitle

\noindent\emph{Source and licence.} Invariant measures with full support and approximation by zero-entropy systems in the $C^0$-Gromov--Hausdorff topology. Version arXiv:2604.02810v1 (CC BY 4.0), distributed by arXiv under the Creative Commons Attribution 4.0 International licence, http://creativecommons.org/licenses/by/4.0/, as recorded in the arXiv licence metadata for this exact version and shown on the abstract page https://arxiv.org/abs/2604.02810v1. Full licence notice: \texttt{licenses/arxiv-2604.02810-CC-BY-4.0.txt}.\par\smallskip

\noindent\textit{Editorial scope.} Counted unit: the article's theorem thm:main (source0.tex lines 346--353). The article's complete original proof of it is delivered with it (source0.tex lines 355--623, one \texttt{proof} environment of 269 lines). No mathematics is written, repaired or re-derived: the statement and the proof are the article's own lines, in the article's own notation, and no retained line is substituted or re-flowed.  Retained uncounted as the source-local context the proof needs: lines 219--224; lines 273--288; lines 292--302.  Nothing was omitted from any retained span, so the package carries no omission record.  No retained line contains a citation, so the wrapper carries no bibliography and no bibliography entry.  No retained line contains a figure environment, an image-inclusion command or drawing code, so no image is delivered and the package declares no asset dependency.  Editorial formatting only, in addition: an AMS \texttt{amsart} preamble that restates the article's own declarations and macros used by the retained lines, namely the environments \texttt{lemma}, \texttt{theorem} on separate counters, together with the standard packages those lines need.  The retained environments print in this document as Lemma 1, Theorem 1 in order of appearance, whereas the article prints the same items with the numbering of its own full document.  The complete native source of the article as distributed in the arXiv e-print for this exact version is bundled unchanged as \texttt{sources/197761/source0.tex}.
\begin{lemma}\label{lem:finite_entropy_zero}
If $Y$ is a finite set and $g:Y\to Y$ is a map, then
\[
h_{\topo}(g)=0.
\]
\end{lemma}

\begin{theorem}[Ergodic decomposition]\label{thm:ergodic_decomposition}
Let $X$ be a compact metric space and let $f:X\to X$ be a homeomorphism. If $\mu$ is an $f$-invariant probability measure, then there exist a measurable set $X_0\subset X$ with $\mu(X_0)=1$, a measurable partition $\mathcal P$ of $X_0$, and a family of probability measures $\{\mu_P\}_{P\in\mathcal P}$ such that:
\begin{itemize}
    \item[(i)] for $\hat{\mu}$-almost every class $P\in\mathcal P$, one has $\mu_P(P)=1$;
    \item[(ii)] for every measurable set $E\subset X$, the map
    \[
    P\mapsto \mu_P(E)
    \]
    is measurable;
    \item[(iii)] for $\hat{\mu}$-almost every class $P\in\mathcal P$, the measure $\mu_P$ is invariant and ergodic;
    \item[(iv)] for every measurable set $E\subset X$,
    \[
    \mu(E)=\int \mu_P(E)\,d\hat{\mu}(P).
    \]
\end{itemize}
\end{theorem}

\begin{lemma}\label{lem:dense_orbit_support}
Let $f:X\to X$ be a continuous map on a compact metric space, and let $\mu$ be an invariant ergodic probability measure. Then there exists a point $\omega \in \supp(\mu)$ such that
\[
\overline{\mathcal O_f^+(\omega)}=\supp(\mu),
\]
where
\[
\mathcal O_f^+(\omega)=\{\omega,f(\omega),f^2(\omega),\dots\}.
\]
In particular, the restriction $f|_{\supp(\mu)}$ is topologically transitive.
\end{lemma}

\begin{theorem}\label{thm:main}
Let $f:X\to X$ be a homeomorphism of a compact metric space. If $f$ admits an invariant probability measure with full support, then for every $\delta>0$ there exist a compact metric space $Y_\delta$ and a homeomorphism $g_\delta:Y_\delta\to Y_\delta$ such that
\[
d_{GH^0}(f,g_\delta)<\delta
\qquad\text{and}\qquad
h_{\topo}(g_\delta)=0.
\]
\end{theorem}

\begin{proof}
Let $\mu$ be an $f$-invariant probability measure with full support. Fix $\delta>0$. Since $f$ is uniformly continuous and $X$ is compact, there exists $\beta>0$ such that
\begin{eqnarray}\label{ecu0}
d(x,y)<\beta
\quad\Longrightarrow\quad
d\bigl(f(x),f(y)\bigr)<\frac{\delta}{3}
\qquad\text{for all }x,y\in X.
\end{eqnarray}
Take $0<\alpha<\min \left\{\frac{\delta}{9},\frac{\beta}{2}\right\}$. By compactness of $X$, there exists a finite set $\{x_1,\dots,x_r\}\subset X$ such that
\[
d_H(\{x_1,\dots,x_r\},X)<\alpha.
\]
Let $\{\mu_P\}_{P\in\mathcal P}$ be the ergodic decomposition of $\mu$ as in Theorem \ref{thm:ergodic_decomposition}. Since $\supp(\mu)=X$, one has
\[
\mu\bigl(B(x_i,\alpha)\bigr)>0
\qquad\text{for every } i=1,\dots,r.
\]
Applying Theorem \ref{thm:ergodic_decomposition}, for each $i\in\{1,\dots,r\}$ we obtain
\[
\mu\bigl(B(x_i,\alpha)\bigr)
=
\int \mu_P\bigl(B(x_i,\alpha)\bigr)\,d\hat\mu(P)>0.
\]
Since $\mu_P$ is ergodic for $\hat\mu$-almost every class $P\in\mathcal P$, there exists a class $P_i\in\mathcal P$ such that
\begin{eqnarray}\label{ecu1}
\mu_{P_i}\bigl(B(x_i,\alpha)\bigr)>0,
\end{eqnarray}
and $\mu_{P_i}$ is ergodic. Denote this measure by $\mu_i$. From \eqref{ecu1}, for every $i=1,\ldots,r$, one has
\begin{eqnarray}\label{ecu2}
B(x_i,\alpha)\cap \supp(\mu_i)\neq\emptyset.
\end{eqnarray}

By Lemma \ref{lem:dense_orbit_support}, for each $i$ there exists a point $z_i\in \supp(\mu_i)$ such that
\[
\overline{\mathcal O_f^+(z_i)}=\supp(\mu_i).
\]
Since $B(x_i,\alpha)\cap \supp(\mu_i)\neq\emptyset$ and the positive orbit of $z_i$ is dense in $\supp(\mu_i)$, there exists an integer $\ell_i\ge0$ such that
\[
f^{\ell_i}(z_i)\in B(x_i,\alpha).
\]
Define
\[
a_i:=f^{\ell_i}(z_i).
\]
Since $f^{\ell_i}$ is a homeomorphism and $\supp(\mu_i)$ is $f$-invariant, one has
\begin{equation}\label{ecu3}
\overline{\mathcal O_f^+(a_i)}
=
\overline{f^{\ell_i}\bigl(\mathcal O_f^+(z_i)\bigr)}
=
f^{\ell_i}\bigl(\overline{\mathcal O_f^+(z_i)}\bigr)
=
f^{\ell_i}\bigl(\supp(\mu_i)\bigr)
=
\supp(\mu_i).
\end{equation}
In particular,
\[
a_i\in \supp(\mu_i)\cap B(x_i,\alpha).
\]
Since $\supp(\mu_i)$ is $f$-invariant and $a_i\in \supp(\mu_i)$, one has
\[
f^{-1}(a_i)\in \supp(\mu_i).
\]
By \eqref{ecu3}, since $f^{-1}(a_i)\in \supp(\mu_i)$, there exists an integer $N_i\ge0$ such that
\[
d\bigl(f^{N_i}(a_i),f^{-1}(a_i)\bigr)<\beta.
\]
Applying $f$ and using the choice of $\beta$ in \eqref{ecu0}, we obtain
\[
d\bigl(f^{N_i+1}(a_i),a_i\bigr)
=
d\bigl(f(f^{N_i}(a_i)),f(f^{-1}(a_i))\bigr)
<
\frac{\delta}{3}.
\]

We now define the finite set
\[
Y_\delta=\{(i,k):\,1\le i\le r,\ 0\le k\le N_i\}.
\]
On $Y_\delta$ we consider the map
\[
q:Y_\delta\to X,
\qquad
q(i,k)=f^k(a_i).
\]

To avoid possible identifications between distinct pairs $(i,k)$ and $(j,\ell)$ having the same image in $X$, we equip $Y_\delta$ with the metric
\[
\rho(u,v)=d\bigl(q(u),q(v)\bigr)+\alpha\,\eta(u,v),
\qquad u,v\in Y_\delta,
\]
where
\[
\eta(u,v)=
\begin{cases}
0,& u=v,\\
1,& u\neq v.
\end{cases}
\]
Since $d(q(u),q(v))$ defines a pseudometric on $Y_\delta$ and $\alpha\,\eta(u,v)$ is a metric, the function $\rho$ is a metric on $Y_\delta$. In particular, $(Y_\delta,\rho)$ is a compact metric space.

We now define $g_\delta:Y_\delta\to Y_\delta$ by
\[
g_\delta(i,k)=(i,k+1), \qquad 0\le k\le N_i-1,
\]
and
\[
g_\delta(i,N_i)=(i,0).
\]
Clearly, $g_\delta$ is a permutation of $Y_\delta$; since $Y_\delta$ is finite, it follows that $g_\delta$ is a homeomorphism. Moreover, for every $u=(i,k)\in Y_\delta$ one has
\begin{eqnarray}\label{ecu1.2}
d\bigl(q(g_\delta(u)),f(q(u))\bigr)<\frac{\delta}{3}.
\end{eqnarray}
Indeed, if $0\le k\le N_i-1$, then
\[
q(g_\delta(i,k))=q(i,k+1)=f^{k+1}(a_i)=f(q(i,k)),
\]
whereas if $k=N_i$, then
\[
d\bigl(q(g_\delta(i,N_i)),f(q(i,N_i))\bigr)
=
d\bigl(a_i,f^{N_i+1}(a_i)\bigr)
<
\frac{\delta}{3}.
\]
On the other hand, since $\{a_1,\dots,a_r\}\subset q(Y_\delta)$, we have
\[
d_H(q(Y_\delta),X)\le d_H(\{a_1,\dots,a_r\},X).
\]
Moreover,
\[
d_H(\{a_1,\dots,a_r\},X)
\le
 d_H(\{a_1,\dots,a_r\},\{x_1,\dots,x_r\})
+
d_H(\{x_1,\dots,x_r\},X)
<
2\alpha,
\]
since $a_i\in B(x_i,\alpha)$ for every $i$. Therefore,
\begin{equation}\label{ecu6}
d_H(q(Y_\delta),X)<2\alpha<\delta.
\end{equation}

Furthermore, if $u\neq v$, then
\[
\rho(u,v)=d\bigl(q(u),q(v)\bigr)+\alpha.
\]
Therefore, for every $u,v\in Y_\delta$ one has $\bigl|d\bigl(q(u),q(v)\bigr)-\rho(u,v)\bigr|\le \alpha$. 
Hence,
\begin{eqnarray}\label{ecu7}
\sup_{u,v\in Y_\delta}
\bigl|d\bigl(q(u),q(v)\bigr)-\rho(u,v)\bigr|
\le \alpha<\delta.
\end{eqnarray}
From \eqref{ecu6} and \eqref{ecu7} we conclude that
\[
q:(Y_\delta,\rho)\to X
\]
is a $\delta$-isometry. Moreover,
\[
d_{C^0}(q\circ g_\delta,f\circ q)<\frac{\delta}{3}<\delta.
\]

Since $d_H(q(Y_\delta),X)<2\alpha$, we construct a map $j:X\to Y_\delta$ as follows: for each $x\in X$, choose a point $j(x)\in Y_\delta$ such that
\begin{eqnarray}\label{ecu6.1}
d\bigl(x,q(j(x))\bigr)<2\alpha.
\end{eqnarray}
We claim that $j$ is a $\delta$-isometry. Given $x,x'\in X$, one has
\[
\bigl|d\bigl(q(j(x)),q(j(x'))\bigr)-d(x,x')\bigr|
\le
 d\bigl(x,q(j(x))\bigr)+d\bigl(x',q(j(x'))\bigr)
<
4\alpha.
\]
Moreover, by the definition of $\rho$,
\[
\bigl|\rho(j(x),j(x'))-d\bigl(q(j(x)),q(j(x'))\bigr)\bigr|
\le \alpha.
\]
Therefore,
\begin{eqnarray}\label{ecu8}
\bigl|\rho(j(x),j(x'))-d(x,x')\bigr|
<
5\alpha
<
\delta,
\end{eqnarray}
since $\alpha<\delta/9$. On the other hand, for every $y\in Y_\delta$,
\[
\rho\bigl(j(q(y)),y\bigr)
\le
 d\bigl(q(j(q(y))),q(y)\bigr)+\alpha
<
3\alpha.
\]
Therefore,
\begin{eqnarray}\label{ecu9}
d_H\bigl(j(X),Y_\delta\bigr)<3\alpha<\delta.
\end{eqnarray}
From \eqref{ecu8} and \eqref{ecu9}, we conclude that $j:X\to (Y_\delta,\rho)$
is a $\delta$-isometry.
\vspace{0.3cm}

We now verify that $d_{C^0}(j\circ f,g_\delta\circ j)<\delta$. Indeed, let $x\in X$. By the definition of $\rho$, one has
\[
\rho\bigl(j(f(x)),g_\delta(j(x))\bigr)
\le
 d\bigl(q(j(f(x))),q(g_\delta(j(x)))\bigr)+\alpha.
\]
Moreover,
\[
\begin{aligned}
d\bigl(q(j(f(x))),q(g_\delta(j(x)))\bigr)
&\le
 d\bigl(q(j(f(x))),f(x)\bigr)
+d\bigl(f(x),f(q(j(x)))\bigr)\\
&\qquad
+d\bigl(f(q(j(x))),q(g_\delta(j(x)))\bigr).
\end{aligned}
\]
By the definition of $j$ given in \eqref{ecu6.1}, the first term on the right-hand side is smaller than $2\alpha$. For the second term, since
\[
d\bigl(x,q(j(x))\bigr)<2\alpha<\beta,
\]
the choice of $\beta$ in \eqref{ecu0} implies that $d\bigl(f(x),f(q(j(x)))\bigr)<\frac{\delta}{3}$.
Finally, by \eqref{ecu1.2}, the third term is smaller than $\frac{\delta}{3}$. Consequently,
\[
\rho\bigl(j(f(x)),g_\delta(j(x))\bigr)
<
2\alpha+\frac{\delta}{3}+\frac{\delta}{3}+\alpha
=
3\alpha+\frac{2\delta}{3}
<
3\left(\frac{\delta}{9}\right)+\frac{2\delta}{3}
=
\delta.
\]
Therefore,
\[
d_{C^0}(j\circ f,g_\delta\circ j)<\delta.
\]

We have constructed $\delta$-isometries
\[
q:(Y_\delta,\rho)\to X
\qquad\text{and}\qquad
j:X\to (Y_\delta,\rho)
\]
such that
\[
d_{C^0}(q\circ g_\delta,f\circ q)<\delta
\qquad\text{and}\qquad
d_{C^0}(j\circ f,g_\delta\circ j)<\delta.
\]
By the definition of the distance $d_{GH^0}$, it follows that
\[
d_{GH^0}(f,g_\delta)<\delta.
\]

Finally, since $Y_\delta$ is finite, Lemma \ref{lem:finite_entropy_zero} implies that
\[
h_{\topo}(g_\delta)=0.
\]
This completes the proof.
\end{proof}

\end{document}
